%% ==========================================================================
%% SECTION 1: INTRODUCTION
%% ==========================================================================
\balance
\section{Introduction}\label{sec:intro}

Systems that incorporate language model components face a verification problem.
The components generate text by sampling from learned probability distributions,
and some fraction of that text is fabricated: fluent, confident, and wrong.
LLM based systems today are rife with ``hallucinations''.

% This paper provides the formal framework and empirical cost surface for making that decision.

We prefer the term \emph{fabrication} throughout this paper to emphasize that
fluent confabulation is the system's default behavior under coherence
optimization, not a malfunction to be patched. A language model optimized for
coherence produces coherent output. When the training data provides grounding,
that output tends to track truth. When it does not, the same optimization
produces fluent, confident fabrication. Xu et
al.~\cite{xu2024hallucination} prove that hallucination is
inevitable for LLMs over the class of computable functions.
%Our result is orthogonal. They prove a computational limit on what the model can produce;
% we prove an observational limit on what a supervisor can verify.
The model may know it is fabricating. The supervisor cannot tell
and must verify the generated output.

However, verification is not free or cheap. Every system that
deploys a language model component makes, explicitly or implicitly, a budget
allocation decision: how much of the output to verify, which outputs to
prioritize, and what signals to use for triage.
For example, checking a citation against a database costs API
calls and latency, having a domain expert review a medical summary costs time
and salary, running a second model as a judge costs compute.
Preliminary observations suggest verification cost varies by content format:
natural text, documented code, compact code, and mathematical proofs exhibit
different scaffolding ratios (syntactic vs. semantic), predicting different optimal judge strategies.

The underlying issue is that current models
expose a text-only output interface. The standard
text-only interface exports only the model's chosen output: its
\emph{testimony} (what the model says).
Under the current training objectives of coherency, the testimony
lies leading to fabrications.

This creates an observability gap with direct consequences for
verification cost. A supervisor receiving only text must spend its
verification budget reconstructing, from the output alone,
what the model's own computation already knew. The budget is
consumed by reconstruction rather than concentrated on the cases
where verification is most needed.

We introduce epistemic observability for LLMs that exposes epistemic
telemetry of the model produced during the generative procedure.
As part of epistemic observability, we introduce two key components, entropy traces
and attention summaries, that encode the epistemic metadata.

The key insight behind epistemic observability is that language models carry internal
signals, such as entropy traces, attention geometry, token-level
log-probabilities, that distinguish grounded generation from
fabrication. These signals are byproducts of the computation itself:
\emph{telemetry} (measurements of the computation that produced the output)
that the model cannot separately control under standard training.
Unlike testimony, the telemetry is harder to fake.
%Yet the interface discards the telemetry and delivers only the testimony.

This paper establishes three contributions that together characterize the cost
surface of epistemic verification for fabricated output.

\fakepara{An impossibility result as budget constraint.}: epistemic honesty is unverifiable for
systems operating under text-only observation with bounded supervision
(\autoref{sec:formal}).
The impossibility is structural: it holds regardless
of model scale, training procedure, or prompt engineering, and
cannot be escaped by stacking supervisors.
This impossibility functions as a budget constraint: below this
investment line, text-channel verification cannot help.

\fakepara{Epistemic Cost Surface.}
Systems engineers building with language
model components need to know what each level of verification investment
buys, and where the investment transitions from text-channel (cheap,
gameable) to tensor-channel (more expensive, robust).
We provide an \emph{honest decomposition}
of verification effectiveness into text-channel and tensor-channel
components, characterize the cost-ro\-bust\-ness tradeoff, and identify
the format-dependent variation that determines judge strategy per domain.
The gap between the text-channel
ceiling and the tensor-channel floor is the impossibility theorem made
empirical: it measures the boundary between gameable and robust verification.

\fakepara{Epistemic Observability-based tensor interface}. We construct a \emph{tensor interface} that escapes the impossibility
class by exporting computational signals alongside text
(\autoref{sec:escape}). The interface requires no architectural changes to
current transformers. It exports entropy traces and attention summaries, which is telemetry
that the model cannot independently control. Under bounded verification,
tensor-guided judges at 10\% budget outperform text-guided judges at 30\%
budget (82.1\% vs.\ 78.5\% accuracy), not because the tensor signal is
strong in isolation, but because it enables efficient triage: concentrating
verification on the cases most likely to contain fabrications.
